\section{Cross-Shard Transaction Model}
\label{sec:cross-shard-model}

\subsection{System abstraction}
\label{subsec:system-abstraction}

We study a cross-shard transaction protocol in which a transfer moves value from
a source shard to a target shard while preserving a consistent terminal outcome.
Each shard maintains an independent UTXO-based ledger, and each cross-shard
transfer is represented by a uniquely identified protocol session.

A session coordinates source locking, cross-shard receipt creation and delivery,
destination preparation, and the final commit decision, while also handling
aborts, validation failures, insufficient quorum, timeouts, duplicate receipts,
and partially applied commits.

The implementation separates shard management from transaction-state mutation.
The shard manager maintains topology, validators, logical time, session
registration, and aggregate metrics, while the atomic-commit component validates
protocol preconditions and applies state changes. This separation allows the
transaction protocol to be analyzed independently from unrelated functionality
in the surrounding simulation platform.

Receipt authenticity is an abstract protocol precondition rather than a
cryptographic result. Quorum summarizes validator approval without modeling the
consensus mechanism that produces votes, and timeouts use logical rounds rather
than wall-clock deadlines.

\subsection{Cross-shard session state machine}
\label{subsec:state-machine}

Each transfer is represented by a finite-state session. The implementation
contains the following states:

\begin{itemize}
\item \texttt{CREATED}
\item \texttt{SOURCE\_\allowbreak LOCKED}
\item \texttt{RECEIPT\_\allowbreak CREATED}
\item \texttt{RECEIPT\_\allowbreak DELIVERED}
\item \texttt{DESTINATION\_\allowbreak PREPARED}
\item \texttt{COMMITTED}
\item \texttt{ABORTED}
\item \texttt{TIMED\_\allowbreak OUT}
\item \texttt{FAILED\_\allowbreak VALIDATION}
\end{itemize}

The normal successful execution follows

\begin{equation}
\begin{split}
\texttt{CREATED}
&\rightarrow \texttt{SOURCE\_\allowbreak LOCKED} \\
&\rightarrow \texttt{RECEIPT\_\allowbreak CREATED} \\
&\rightarrow \texttt{RECEIPT\_\allowbreak DELIVERED} \\
&\rightarrow \texttt{DESTINATION\_\allowbreak PREPARED} \\
&\rightarrow \texttt{COMMITTED}.
\end{split}
\label{eq:normal-commit-path}
\end{equation}

The states \texttt{COMMITTED}, \texttt{ABORTED}, \texttt{TIMED\_\allowbreak OUT},
and \texttt{FAILED\_\allowbreak VALIDATION} are terminal and have no outgoing
transitions.

Abort and validation-failure transitions may occur from nonterminal states. A
timeout may occur after source locking and before a terminal decision, but not
directly from \texttt{CREATED}, where no source resource has been locked.

The transition relation rejects commit after abort, commit after timeout,
creation-to-commit jumps, repeated terminal decisions, and receipt processing
for a different transfer identifier. Each accepted transition emits an
immutable event containing sequence number, logical time, transfer identifier,
action, previous state, next state, and diagnostic reason.

\subsection{Atomic commit and rollback}
\label{subsec:atomic-commit}

The executable commit operation has four conceptual stages:

\begin{enumerate}
\item preparation of the commit plan
\item validation of commit preconditions
\item application of ledger and session mutations
\item rollback if an application step fails
\end{enumerate}

Before modifying ledger state, the protocol constructs a commit plan and
captures a restorable snapshot containing the source UTXO and lock state,
receipt-consumption state, optional change output, destination output, and a
checkpoint of the complete protocol session.

During normal commit, the destination is prepared, the receipt consumed, the
source value removed, the source lock released, optional change created, the
destination output created, and the session moved to \texttt{COMMITTED}.

A failure invokes rollback before the protocol error is propagated. Rollback
restores the destination and change outputs, source UTXO, receipt-consumption
state, source lock, and session checkpoint. No synthetic reverse transition is
recorded because rollback cancels an incomplete state application rather than
creating a new protocol decision.

Deterministic failure-injection points around receipt consumption, source debit,
and destination credit make partially executed commits reproducible without
introducing nondeterministic faults.

\subsection{Concurrent formal abstraction}
\label{subsec:formal-abstraction}

Bounded TLA+ and Alloy models abstract the executable protocol to the
protocol-level state required for concurrent sessions and failure behavior. They
represent transfer status, source and target shards, resource locks, receipt
ownership and consumption, destination credit, released funds, messages, and
validator votes, while collapsing implementation states that do not alter the
protocol-level decision state.

The Java-to-formal projection maps \texttt{CREATED} to \texttt{Pending};
\texttt{SOURCE\_\allowbreak LOCKED}, \texttt{RECEIPT\_\allowbreak CREATED}, and
\texttt{RECEIPT\_\allowbreak DELIVERED} to \texttt{Locked};
\texttt{DESTINATION\_\allowbreak PREPARED} to \texttt{Prepared};
\texttt{COMMITTED} to \texttt{Committed}; and \texttt{ABORTED},
\texttt{TIMED\_\allowbreak OUT}, and \texttt{FAILED\_\allowbreak VALIDATION} to
\texttt{Aborted}.

TLA+ additionally stores the first abstract terminal decision in
\texttt{terminalStatus}. Alloy instead represents bounded executions as an
ordered sequence of states and checks terminal-state irreversibility over
consecutive states. The two models therefore encode terminal consistency
differently while targeting the same protocol concern.

Concurrent configurations vary shards and active transfers and include
duplicate receipts, delayed messages, timeout races, and quorum behavior. They
are finite by construction, so conclusions are bounded by the declared scopes
and constants.

TLA+ and Alloy are complementary representations of the same protocol concerns,
not equivalent state spaces or interchangeable cost models. Absolute
execution-time differences are therefore not used to infer tool superiority.

\subsection{Evaluated protocol properties}
\label{subsec:properties}

Seven protocol properties define the main bounded verification target.

\begin{description}

\item[] \texttt{NoReceipt\allowbreak Replay}.
A cross-shard receipt must not be consumed more than once.

\item[] \texttt{DestinationCreditRequires\allowbreak ValidReceipt}.
Destination-side credit requires a receipt accepted as valid by the protocol.

\item[] \texttt{Decision\allowbreak Consistency}.
A committed transfer must not simultaneously be represented as having released
its source funds through an abort-like terminal outcome.

\item[] \texttt{NoValueLossAt\allowbreak Termination}.
At a terminal state, transferred value must either have been committed to the
destination or remain recoverable on the source side.

\item[] \texttt{TerminalState\allowbreak Irreversibility}.
Once the protocol reaches its first terminal decision, a later execution step
must not replace it with an incompatible terminal outcome.

\item[] \texttt{EventuallyReleasedAfter\allowbreak Timeout}.
An abstract transfer represented as \texttt{Aborted} must have released its
source funds. This property captures the release obligation associated with the
modeled timeout and non-commit terminal behavior.

\item[] \texttt{Quorum\allowbreak Required}.
A commit requires the relevant validator quorum. TLA+ parameterizes this
threshold through \texttt{Quorum}, whereas the evaluated Alloy model requires at
least two validator votes.

\end{description}

Together these properties cover replay protection, destination authorization,
decision consistency, value conservation, terminal-state stability, timeout
release, and quorum enforcement. The same set is evaluated for valid formal
models and the mutation-based experiments.

\subsection{Scope of the model}
\label{subsec:model-scope}

Receipt authenticity remains abstract, validator quorum is an approval
condition rather than a verified consensus protocol, and time is represented by
logical rounds. Network behavior is scenario-bounded: delay and duplicate
receipts are represented in evaluated configurations, while concrete network
events not structurally represented by the abstraction are projected as
stuttering steps during trace conformance.

Both formal representations operate over finite configurations. Absence of a
counterexample therefore means only that no violation was found for the
evaluated model, property, configuration, and bound, not an unbounded proof of
protocol correctness.

The Java implementation is related to the formal model through bounded trace
conformance, not a general refinement proof or behavioral equivalence.
